\section{A convex potential from inverse-distance kernels}\label{lk:section}

Explicit revelations are one way to pay for a concentrated posterior.
This section uses a different accounting method. From the relative
entropy of a posterior we subtract the logarithm of an inverse-distance
average. The resulting functional is convex, so after a block we may
forget the entering state without increasing it. A mixed projection
estimate then charges only the rows that the learner actually saw.

Throughout this section, \(\sigma\) is uniform probability on
\(S^{d-1}\), \(n=d-1\), \(k=\lfloor d/16\rfloor\), and \(p=2k\).
We work in sufficiently large dimension that \(k\ge1\) and \(2p<n\).
For a probability measure \(\nu\le H\sigma\), \(H<\infty\), put
\begin{equation}\label{lk:definition}
 J_\nu(s)=\int\|u-s\|^{-p}\,d\nu(u),\qquad
 B(\nu)=\int\log J_\nu(s)\,d\nu(s),\qquad
 \Phi(\nu)=D(\nu\|\sigma)-B(\nu).
\end{equation}
The sphere has diameter two, so \(J_\nu\ge2^{-p}\). The centered
ball bound of Lemma~\ref{l:caps-endpoint}, summed on dyadic annuli,
bounds \(J_\nu\) above uniformly in \(s\) because \(p<n\).
Thus every term in \eqref{lk:definition} is finite.

\Needspace{6\baselineskip}
\begin{lemma}[Convexity and coercivity]\label{lk:coercive}
The functional \(\Phi\) is convex under finite mixtures of
bounded-density probability measures and
\[
 \Phi(\nu)\ge\tfrac12D(\nu\|\sigma)-Cd,\qquad
 \Phi(\sigma)\le p\log2
\]
for an absolute \(C\).
\end{lemma}
\begin{proof}
Writing \(f=d\nu/d\sigma\), we have
\(\Phi(\nu)=\int f\log(f/J_\nu)\,d\sigma\), with value zero where
\(f=0\). Both \(f\) and \(J_\nu\) are linear in \(\nu\).
For \(x>0,y>0\), the Hessian of \(x\log(x/y)\) has quadratic form
\((a-bx/y)^2/x\); continuity at \(x=0\) gives joint convexity on
\(x\ge0,y>0\). Integration proves the mixture assertion.

By Jensen and rotational invariance,
\[
 \int J_\nu(s)^2\,d\sigma(s)
 \le\int\!\int\|u-s\|^{-2p}\,d\sigma(s)d\nu(u)
 \le e^{Cd}.
\]
For the last bound, distances greater than one contribute at most
one. On \(2^{-i-1}<\|u-s\|\le2^{-i}\), the contribution is at most
\(2^{2p}2^{-i(n-2p)}\); its sum converges since \(n-2p>0\).
Apply the entropy inequality \eqref{l:tilt}, equivalently
\(\E_\nu F\le D(\nu\|\sigma)+\log\E_\sigma e^F\), to
\(F=2\log J_\nu\). This logarithm is bounded for fixed \(\nu\).
It gives \(2B(\nu)\le D(\nu\|\sigma)+Cd\). Finally,
\(D(\sigma\|\sigma)=0\) and \(J_\sigma\ge2^{-p}\).
\end{proof}

\subsection{Exact projection densities and the kernel entropy}

We will use differential entropy of exact labels. The following
elementary density statement fixes both its version and its
integrability before the information identity is used.

\begin{lemma}[Projection versions and logarithms]\label{lk:density}
Let \(1\le\ell<d-1\), let \(G\) have \(\ell\) independent standard
Gaussian rows, and let \(\nu\le H\sigma\) be a probability measure.
For each full-row-rank \(g\), the image \(g_\#\nu\) has a Lebesgue
density. The limit, where finite and existent, of
\[
 p_{\nu,a}(g,y)=(2a)^{-\ell}
     \nu\{u:\|gu-y\|_\infty\le a\},
 \qquad a=1,1/2,1/3,\ldots,
\]
and zero elsewhere is a jointly Borel density version \(p_\nu(g,y)\).
Under the independent law of \(G\) and \(S\sim\nu\),
\(\log p_\nu(G,GS)\) is integrable.

If instead \(G\) and \(S\sim\sigma\) are marginally independent and
\(W\) is finite with any joint law with them, the actual conditional
law of \(GS\) given \((G,W)\) has a Lebesgue density. Its logarithm
and \(\log p_{\nu_W}(G,GS)\), where
\(\nu_w=P_{S\mid W=w}\), are integrable under the actual law.
The same assertion holds after first conditioning on a finite entering
variable \(V\), provided the fresh \(G\) is jointly independent of
\((S,V)\). Equivalently, conditional on each positive-probability value
of \(V\), the joint law of \((G,S)\) must be the product of the original
Gaussian row law and the conditional law of \(S\).
\end{lemma}
\begin{proof}
Writing \(H_g=(gg^{\mathsf T})^{1/2}\), an orthogonal change of
coordinates and a nonsingular change of label coordinates give
\[
 p_\sigma(g,y)=
 \frac{c_{d,\ell}}{\det(gg^{\mathsf T})^{1/2}}
 \left(1-\|H_g^{-1}y\|^2\right)^{(d-\ell-2)/2}
 \ind_{\{\|H_g^{-1}y\|<1\}}.
\]
The spherical coordinate formula follows by splitting a normalized
standard Gaussian into its \(\ell\) row-space coordinates and the
remaining coordinates: the squared projected length is
\(\operatorname{Beta}(\ell/2,(d-\ell)/2)\). Domination proves absolute
continuity for \(\nu\). Lebesgue differentiation proves the claim
about cube limits, and their defining integrals are jointly Borel.

Under \(G\otimes\sigma\), the logarithm in the displayed formula is
integrable. The beta density has integrable logarithms at its two
endpoints. Gram--Schmidt expresses \(\det(GG^{\mathsf T})\) as a
product of chi-square variables of positive degrees, each with
integrable logarithm. Domination by \(H\) preserves this baseline
integrability for \(S\sim\nu\). The projected likelihood ratio
\(p_\nu/p_\sigma\) has relative entropy at most
\(D(\nu\|\sigma)\le\log H\) by data processing. Its negative
logarithmic part has expectation at most \(1/e\), since
\(t\log^-(t)\le1/e\) under the probability reference. Its positive
part is then finite as well. This proves the first integrability
assertion.

For a positive-probability \(w\), the actual law of \((G,S)\) given
\(w\) is dominated by \(P(W=w)^{-1}\) times \(G\otimes\sigma\).
Its image therefore has a density, and the baseline logarithm is
integrable; averaging over \(w\) restores its original finite
expectation. The log ratio of the actual image density to
\(p_\sigma\) has relative entropy \(I(W;GS\mid G)\le H(W)\) and
negative part at most \(1/e\). Thus the actual log density is
integrable. Moreover,
\[
 \E D(P_{S\mid G,W}\|\nu_W)=I(S;G\mid W)\le H(W),
\]
where the last inequality uses \(I(S;G)=0\) and the chain rule.
Data processing under \(s\mapsto Gs\) gives integrability of the
log ratio between the actual image density and \(p_{\nu_W}\);
subtracting it proves the remaining claim. If a finite entering
variable is first fixed, its signal law is bounded relative to
\(\sigma\), and the same argument applies to each positive-probability
case and then to their finite average.
\end{proof}

\begin{lemma}[The kernel controls projection entropy]\label{lk:entropy}
Let \(X\) have \(p\) independent standard Gaussian rows and be
independent of \(S\sim\nu\le H\sigma\), \(H\ge1\). For a sufficiently
large absolute \(C_0\), put
\(b_H=\lceil C_0(d+\log(2H))\rceil\). Then
\[
 h(XS\mid X)\le -B(\nu)+Cd+2\log(b_H+2).
\]
\end{lemma}
\begin{proof}
For \(r_i=2^{1-i}\), let \(m_i(s)=\nu(B(s,r_i))\).
Dyadic shells give
\[
 J_\nu(s)\le 2^p\sum_{i\ge0}r_i^{-p}m_i(s).
\]
Since \(m_i(s)\le H(2r_i)^n\) and \(n-p\ge1\), choosing \(C_0\)
large makes the sum over \(i>b_H\) at most \(2^{-p}\).
The term \(i=0\) is itself \(2^{-p}\). Choose the first maximizer
\(j=j(s)\le b_H\). This is measurable and satisfies
\[
 \log(r_j^{-p}m_j(s))
 \ge\log J_\nu(s)-Cd-\log(b_H+2).
\]

For fixed \(X\), compare the actual projected density to
\[
 Q_X(y)=\frac1{b_H+1}\sum_{i=0}^{b_H}
       \int \gamma_{r_i}(y-Xu)\,d\nu(u),
\]
where \(\gamma_r\) is the \(N(0,r^2I_p)\) density. At \(y=Xs\),
retain the \(j\)-summand and restrict it to \(B(s,r_j)\).
Jensen's inequality for the logarithm on that normalized restriction
and \(\E_X\|X(s-u)\|^2=p\|s-u\|^2\le pr_j^2\) give
\[
 \E_X\log Q_X(Xs)
 \ge\log(r_j^{-p}m_j(s))-\log(b_H+1)
     -\tfrac p2\log(2\pi)-\tfrac p2.
\]
Integrate over \(s\). The comparison density has finite cross
entropy by this bound and its finite Gaussian-mixture upper bound.
The actual differential entropy is finite by Lemma~\ref{lk:density}.
Nonnegativity of relative entropy to \(Q_X(y)\,dy\) proves the result.
\end{proof}

\subsection{The mixed estimate and the actual-row charge}

The following is the positive-scale assertion of
\cite[Theorem~3.1]{OpenAIRow2026}. We state its parameters explicitly
because its application has only one unit of angular slack. Its proof
is in that companion; the reweighting and change of row law below are
proved here.

\begin{theorem}[Mixed projection estimate]\label{lk:mixed-input}
Let \(m,r,q\ge1\) be integers, \(\ell=m+r<n=d-1\), and
\(\ell-m-(q-1)\ge1\). Let \(\rho\) be a finite measure on the sphere
such that, for every ambient \(z\) and \(t>0\),
\(\rho(B(z,t))\le Bt^\ell\) and
\(\rho(B(z,t))\le Dt^n\), with \(B,D>0\). For independent Gaussian
matrices \(X,Z\) with \(m,r\) rows, put \(G=(X;Z)\) and
\[
 p_{\rho,\delta}(G,y)=(2\delta)^{-\ell}
      \int\ind_{\{\|Gu-y\|_\infty\le\delta\}}\,d\rho(u).
\]
For every fixed \(s\in S^{d-1}\) and every \(\delta>0\),
\[
 \left[\E_X\left(\E_Zp_{\rho,\delta}(G,Gs)\right)^q\right]^{1/q}
 \le e^{Cd}B\bigl(1+\log^+(D/B)\bigr).
\]
\end{theorem}

The estimate is for positive cube averages and every fixed signal.
In its use below, we first change the law of the actual rows while
\(\delta>0\), and only then pass to exact labels under their actual
law. No assertion about a fixed-signal value of an arbitrary density
version is used.

\begin{lemma}[Reweighting by the potential]\label{lk:reweight}
For \(\nu\le H\sigma\), define \(d\rho=d\nu/J_\nu\). Then
\[
 \rho(B(z,r))\le(2r)^p,\qquad \rho\le2^pH\sigma.
\]
By Theorem~\ref{lk:mixed-input}, for independent \(k\)-row
Gaussian matrices \(A,A'\), \(X=(A;A')\), every \(s\in\sphere\), and
every \(\delta>0\),
\begin{equation}\label{lk:moment}
 \left[\E_A\left(\E_{A'}p_{\rho,\delta}(X,Xs)\right)^k\right]^{1/k}
 \le e^{Cd}(b_H+2).
\end{equation}
\end{lemma}
\begin{proof}
For \(s\in B(z,r)\), \(J_\nu(s)\ge(2r)^{-p}\nu(B(z,r))\).
Integrating its reciprocal over the ball proves the first inequality;
if the ball has zero \(\nu\)-mass there is nothing to prove. The second
uses \(J_\nu\ge2^{-p}\).

Apply Theorem~\ref{lk:mixed-input} with
\(m=r=q=k\), \(\ell=p=2k\), \(B=2^p\), and
\(D=2^{p+n}H\), using the ambient-center cap bound of
Lemma~\ref{l:caps-endpoint}. The angular margin is exactly
\(p-k-(k-1)=1\), and the radial margin is \(n-p>0\).
Moreover \(1+\log^+(D/B)\le C(d+1+\log H)\le C(b_H+2)\);
the factor \(2^p\) is absorbed in \(e^{Cd}\). This proves
\eqref{lk:moment}. The local inverse factor behind this substitution
is \(\|u-s\|^{-k}\dist(u-s,E)^{-k}\), where \(E\) is a linear
subspace with \(\dim E\le k-1\).
\end{proof}

\begin{proposition}[Drift of the kernel functional]\label{lk:drift}
Let \(S\sim\sigma\), let
\(V\) have at most \(N\) values, and let a \(k\)-row Gaussian
matrix \(A\) be independent of \((S,V)\). Suppose \(W\) has at most
\(N\) values and is generated by a Borel probability kernel of
\((V,A,AS)\) and independent randomness. If \(\log N\le d^2\),
then, for sufficiently large \(d\),
\[
 \E\Phi(P_{S\mid W})\le \E\Phi(P_{S\mid V})+Cd.
\]
\end{proposition}
\begin{proof}
Condition first on \(V=v\) with positive probability. Put
\(\nu=P_{S\mid v}\) and \(H_v=P(V=v)^{-1}\). In this conditional
experiment \(w\) always ranges over values with \(P(V=v,W=w)>0\);
for these values put
\(\nu_w=P_{S\mid v,w}\) and \(H_{vw}=P(V=v,W=w)^{-1}\).
The two posteriors are bounded by \(H_v\sigma\) and
\(H_{vw}\sigma\). Introduce a \(k\)-row Gaussian matrix \(A'\)
independent of \((S,V,A,W)\), and put \(X=(A;A')\), \(Y=XS\).
Within the conditional experiment \(X\) is independent of \(S\),
and \(W\) depends on \(S\) only through \((X,Y)\). The chain rule gives
\begin{equation}\label{lk:identity}
 I(S;W)=h(Y\mid X)-h(Y\mid X,W)-I(X;S\mid W).
\end{equation}
For example,
\(I(S;W)=I(X,Y;W)-I(X;W\mid S)\), and \(I(X;S)=0\)
converts the difference of the two row-information terms into
\(-I(X;S\mid W)\). All displayed differential entropies are finite
by Lemma~\ref{lk:density}, and \(I(X;S\mid W)\le H(W)\).

For fixed \(w\), use the finite nonzero measure
\(\rho_w=\nu_w/J_{\nu_w}\) as a reference, and write
\(a_w=\rho_w(\sphere)\). First normalize this reference.
The relative entropy of the actual conditional law of \((X,S)\)
with respect to \(P_{X\mid w}\otimes(\rho_w/a_w)\) is
\[
 I(X;S\mid w)+B(\nu_w)+\log a_w<\infty.
\]
Indeed, \(\log J_{\nu_w}\) is bounded for this \(v,w\), and
\(I(X;S\mid w)<\infty\). Probability data processing under
\((X,s)\mapsto(X,Xs)\) makes the corresponding post-map relative
entropy finite. For probability laws \(P\ll Q\), the elementary bound
\[
 \int_{\{dP/dQ<1\}}-\log(dP/dQ)\,dP\le 1/e
\]
shows that a finite relative entropy has an integrable log density
ratio. The post-map ratio is the actual conditional label density
divided by \(p_{\rho_w}(X,Y)/a_w\). Its logarithm is therefore
integrable. Lemma~\ref{lk:density} makes the logarithm of the actual
conditional label density integrable as well, so
\(\log p_{\rho_w}(X,Y)\) is integrable under the actual conditional
law. Here \(p_{\rho_w}\) is the cube-limit density version; the same
construction applies to a finite measure.

We may now expand both finite-reference divergences and subtract
\(\log a_w\) from probability data processing. Before the map the
divergence is \(I(X;S\mid w)+B(\nu_w)\), while after the map it is
\(-h(Y\mid X,w)-\E[\log p_{\rho_w}(X,Y)\mid w]\). Consequently,
\begin{equation}\label{lk:finite-reference}
 -h(Y\mid X,w)-I(X;S\mid w)
 \le B(\nu_w)+\E[\log p_{\rho_w}(X,Y)\mid w].
\end{equation}
This is precisely the finite-reference data-processing statement
proved in Section~\ref{l:model}, including when \(\rho_w\) has mass
greater than one.

We now bound the last logarithm under the selected actual rows. Let
\(Q\) be the original Gaussian law of \(A\), still conditional on
\(v\), and write \(\kappa_w(a,as)\) for the Borel block probability
of message \(w\). For \(r_w(s)=\int\kappa_w(a,as)\,Q(da)>0\), choose
the conditional row law
\[
 Q_{s,w}(da)=\frac{\kappa_w(a,as)}{r_w(s)}\,Q(da);
\]
when \(r_w(s)=0\), assign \(Q_{s,w}=Q\). The normalizer is positive
for almost every \((s,w)\) under the actual law, and this is a
conditional version of the row law there. In particular
\(Q_{s,w}\ll Q\). Also
\[
 \E_{s,w}D(Q_{s,w}\|Q)=I(A;W\mid S)\le H(W)\le\log N,
\]
so the divergence is finite for almost every \((s,w)\) under that law.
The fresh \(A'\) retains its independent Gaussian law after
conditioning on \((A,s,w)\). For almost every such pair and for
\(\delta>0\), Jensen in \(A'\), followed by \eqref{l:tilt}
with exponent \(k\), gives
\begin{align}
 &\E[\log(1+p_{\rho_w,\delta}(X,Xs))\mid s,w]\notag\\
 &\quad\le\frac{D(Q_{s,w}\|Q)}k+
 \frac1k\log\E_{A\sim Q}
       \left(1+\E_{A'}p_{\rho_w,\delta}(X,Xs)\right)^k\notag\\
 &\quad\le\frac{D(Q_{s,w}\|Q)}k+Cd+\log(b_{H_{vw}}+2).
 \label{lk:row-charge}
\end{align}
The last step uses the \(L^k\) triangle inequality and
\eqref{lk:moment}, whose fresh moment holds for every fixed
\(s\in\sphere\). The displayed row-charge inequality is used only
for almost every \((s,w)\) under the actual law, while \(\delta\)
is positive.

Here is the exact null-set transfer in that passage. For each
full-rank \(X\), Lebesgue differentiation makes
\(p_{\rho_w,1/a}(X,y)\) converge to \(p_{\rho_w}(X,y)\)
for Lebesgue-almost every \(y\). The failure set is jointly Borel
by the cube definitions. Under the actual conditional law given
\((v,w)\), \((X,S)\) is dominated by a finite multiple of
\(\gamma_p\otimes\sigma\); pushing forward shows that \((X,Y)\)
is absolutely continuous with respect to \(\gamma_p(dX)\,dy\).
Consequently it avoids this failure set, including after the actual
rows have been tilted by \(w\). Fatou applied to the nonnegative
\(\log(1+p_{\rho_w,1/a}(X,Y))\), after integrating \eqref{lk:row-charge}
over \(s,w\), proves the same integrated upper bound for
\(\log(1+p_{\rho_w}(X,Y))\), and hence for \(\log p_{\rho_w}\).
This argument does not invoke an unsmoothed moment at
\(\rho_w\)-almost every fixed signal.

The preceding estimate supplies the quantitative positive-log bound;
integrability of both signs was established before
\eqref{lk:finite-reference}. Substituting
Lemma~\ref{lk:entropy} and
\eqref{lk:finite-reference}--\eqref{lk:row-charge} into
\eqref{lk:identity}, and using
\(I(S;W)=\E_wD(\nu_w\|\sigma)-D(\nu\|\sigma)\), yields
\[
 \E_w\Phi(\nu_w)\le\Phi(\nu)+Cd+2\log(b_{H_v}+2)
       +\E_w\log(b_{H_{vw}}+2)+\frac{\log N}{k}.
\]
Restore the average over \(v\).
\(\E\log H_V=H(V)\le\log N\) and
\(\E\log H_{VW}=H(V,W)\le2\log N\). Jensen bounds the averaged
cutoff terms by \(C\log(d+1+2\log N)\), while
\(\log N/k\le Cd\). For fixed \(w\), \(P_{S\mid W=w}\) is the
finite mixture of \(\nu_w=P_{S\mid v,w}\) over \(v\).
Lemma~\ref{lk:coercive} permits forgetting \(v\), proving the drift.
\end{proof}

\subsection{Precision and the global stopping index}

\begin{corollary}\label{lk:regression}
Let \(M(d)=o(d^2)\) and
\(0<\epsilon(d)\le1/10\). A learner in the model of
Section~\ref{l:model} with uniform-sphere average angular success at
least \(2/3\) requires
\(T(d)\ge c d\log(1/\epsilon(d))\) for an absolute \(c>0\)
in sufficiently large dimension. The threshold may depend on the
memory sequence.
\end{corollary}
\begin{proof}
Fix the shared rule randomness with uniform-prior success at least
\(2/3\), and apply Lemmas~\ref{l:prior-ac} and~\ref{l:borel}.
Keep the internal transition and output randomness. Put
\(L=\log(1/\epsilon)\). Under the product of the marginal laws of
\((S,\widehat S)\), angular success has probability at most
\(\epsilon^{d-1}\), by Lemma~\ref{l:caps-endpoint}.
Equation~\eqref{l:binary} therefore gives
\begin{equation}\label{lk:output}
 I(S;\widehat S)\ge\tfrac23(d-1)L-\log2.
\end{equation}

Assume \(T<dL\). The other case is already sufficient. The terminal
state and stopping index have at most \((T+1)2^M\) values, so the
left side of \eqref{lk:output} is at most
\(M\log2+\log(T+1)\). Since
\(\log(T+2)\le C+\log d+\log(L+1)\le C+\log d+L\),
absorption gives
\[
 L\le C\frac{M+1+\log d}{d},\qquad T=O(d+M).
\]
For \(M=o(d^2)\), pad stopped states with their original terminal
labels and indices and add unused rows to complete the last block.
Every padded layer has at most \(N=(T+2)2^M\) states and
\(\log N=o(d^2)\).

The initial conditional signal law is \(\sigma\), because the model
makes initialization independent of the signal conditional on the
fixed shared rule randomness. Iterate
Proposition~\ref{lk:drift} for \(\lceil T/k\rceil\) blocks and apply
Lemma~\ref{lk:coercive}. For the final padded state \(F\),
\[
 I(S;F)=\E D(P_{S\mid F}\|\sigma)
 \le Cd\bigl(1+\lceil T/k\rceil\bigr).
\]
The output is generated from \(F\) and fresh randomness. Data
processing with \eqref{lk:output} gives
\(L\le C(1+T/d)\). The \(2/3\) instance of
Lemma~\ref{l:caps-endpoint} gives \(T>d/4\) for this same prior
and exact observation model. Thus \(1+T/d\le5T/d\), proving the
claimed sample bound. The global-clock calculation controls only
the padded width; the precision-dependent conclusion comes from
the block drift.
\end{proof}
